\documentclass[11pt,a4paper]{article}
% ==============================================================
%  Préambule commun aux fiches -- Advanced Probabilities (M1)
%  Usage, dans chaque fiche :
%     \documentclass[11pt,a4paper]{article}
%     % ==============================================================
%  Préambule commun aux fiches -- Advanced Probabilities (M1)
%  Usage, dans chaque fiche :
%     \documentclass[11pt,a4paper]{article}
%     % ==============================================================
%  Préambule commun aux fiches -- Advanced Probabilities (M1)
%  Usage, dans chaque fiche :
%     \documentclass[11pt,a4paper]{article}
%     \input{../preambule.tex}
%     \fiche{C5}{Vecteurs gaussiens}{Gaussian vectors}
% ==============================================================
\usepackage[T1]{fontenc}
\usepackage[utf8]{inputenc}
\usepackage[french]{babel}
\usepackage{lmodern}
\usepackage[a4paper,margin=1.9cm,top=1.7cm,bottom=1.9cm]{geometry}
\usepackage{amsmath,amssymb,amsthm,mathtools}
\usepackage[inline]{enumitem}
\usepackage{booktabs,array,tabularx}
\usepackage{xcolor}
\usepackage[most]{tcolorbox}
\usepackage{fancyhdr}
\usepackage{hyperref}
\hypersetup{colorlinks=true,linkcolor=blue!50!black,urlcolor=blue!50!black}

\setlist{itemsep=1pt,topsep=3pt,parsep=0pt}
\setlength{\parindent}{0pt}
\setlength{\parskip}{3pt}
\emergencystretch=3em

% ---------- Identification de la fiche -------------------------
\newcommand{\ficheid}{}
\newcommand{\fichetitre}{}
\newcommand{\fiche}[3]{%
  \renewcommand{\ficheid}{#1}\renewcommand{\fichetitre}{#2}%
  \pagestyle{fancy}\fancyhf{}%
  \fancyhead[L]{\small\textsf{Advanced Probabilities -- Fiche #1}}%
  \fancyhead[R]{\small\textsf{#2}}%
  \fancyfoot[C]{\small\thepage}%
  \begin{center}
    {\LARGE\bfseries Fiche #1 -- #2}\par\smallskip
    {\large\itshape (#3)}\par\smallskip
    {\small Cours de V.\,M.~Lambert, M1 -- semestre 1}
  \end{center}\vspace{-2pt}\hrule\vspace{6pt}}

% ---------- Environnements numérotés (un seul compteur) --------
\theoremstyle{definition}
\newtheorem{definition}{Définition}
\renewcommand{\thedefinition}{\ficheid.\arabic{definition}}
\newtheorem{theoreme}[definition]{Théorème}
\newtheorem{proposition}[definition]{Proposition}
\newtheorem{lemme}[definition]{Lemme}
\newtheorem{corollaire}[definition]{Corollaire}
\newtheorem{exemple}[definition]{Exemple}
\theoremstyle{remark}
\newtheorem{remarque}[definition]{Remarque}

% Résultat hors cours (standard, signalé comme tel)
\theoremstyle{definition}
\newtheorem{horscours}[definition]{Hors cours}

% Idée de preuve
\newenvironment{idee}{\par\small\textit{Idée de preuve.}~}{\par}

% ---------- Encadrés -----------------------------------------
\tcbset{fichebox/.style={enhanced,breakable,boxrule=0.6pt,arc=2pt,
  left=5pt,right=5pt,top=3pt,bottom=3pt,fonttitle=\bfseries}}
\newtcolorbox{piege}[1][]{fichebox,colback=red!4,colframe=red!60!black,
  title={Pièges classiques},#1}
\newtcolorbox{liens}[1][]{fichebox,colback=blue!3,colframe=blue!50!black,
  title={Liens entre notions},#1}
\newtcolorbox{methode}[1][Méthode]{fichebox,colback=green!4,colframe=green!40!black,title={#1}}
\newtcolorbox{alerte}[1][]{fichebox,colback=orange!6,colframe=orange!70!black,
  title={Attention à l'énoncé},#1}

% Exercice de TD (fiches de rappels)
\newcounter{exo}
\newenvironment{exo}[1][]{\refstepcounter{exo}\par\medskip
  \noindent{\large\bfseries Exercice~\theexo}\ifx\relax#1\relax\else~{\large\bfseries -- #1}\fi\par\nopagebreak\smallskip}{\par}
\newcommand{\outils}{\textbf{Outils du cours.}~}
\newcommand{\meth}{\textbf{Méthode / astuce.}~}
\newcommand{\indic}{\textbf{Indication.}~}
\newcommand{\rep}{\textbf{Réponse.}~}

% ---------- Notations du cours ---------------------------------
\newcommand{\R}{\mathbb{R}}
\newcommand{\N}{\mathbb{N}}
\newcommand{\Q}{\mathbb{Q}}
\newcommand{\C}{\mathbb{C}}
\newcommand{\PP}{\mathbb{P}}
\newcommand{\QQ}{\mathbb{Q}}
\newcommand{\E}{\mathbb{E}}
\newcommand{\F}{\mathcal{F}}
\newcommand{\G}{\mathcal{G}}
\newcommand{\Hc}{\mathcal{H}}
\newcommand{\Tc}{\mathcal{T}}
\newcommand{\Bc}{\mathcal{B}}
\newcommand{\Cc}{\mathcal{C}}
\newcommand{\Ec}{\mathcal{E}}
\newcommand{\Pc}{\mathcal{P}}
\newcommand{\Nc}{\mathcal{N}}
\newcommand{\LN}{\mathcal{LN}}
\newcommand{\Expo}{\mathcal{E}}
\newcommand{\Poi}{\mathcal{P}}
\newcommand{\Ff}{\mathbb{F}}
\newcommand{\1}{\mathbf{1}}
\newcommand{\dd}{\,\mathrm{d}}
\DeclareMathOperator{\Var}{\mathbb{V}}
\DeclareMathOperator{\Cov}{Cov}
\DeclareMathOperator{\card}{card}
\DeclareMathOperator{\diag}{diag}
\newcommand{\indep}{\perp\!\!\!\perp}
\newcommand{\ps}{\ \PP\text{-p.s.}}
\newcommand{\tendloi}{\xrightarrow[n\to+\infty]{\text{loi}}}
\newcommand{\tendps}{\xrightarrow[n\to+\infty]{\text{p.s.}}}
\newcommand{\tendproba}{\xrightarrow[n\to+\infty]{\PP}}
\newcommand{\tendLp}{\xrightarrow[n\to+\infty]{L^p}}
\newcommand{\tend}{\xrightarrow[n\to+\infty]{}}

%     \fiche{C5}{Vecteurs gaussiens}{Gaussian vectors}
% ==============================================================
\usepackage[T1]{fontenc}
\usepackage[utf8]{inputenc}
\usepackage[french]{babel}
\usepackage{lmodern}
\usepackage[a4paper,margin=1.9cm,top=1.7cm,bottom=1.9cm]{geometry}
\usepackage{amsmath,amssymb,amsthm,mathtools}
\usepackage[inline]{enumitem}
\usepackage{booktabs,array,tabularx}
\usepackage{xcolor}
\usepackage[most]{tcolorbox}
\usepackage{fancyhdr}
\usepackage{hyperref}
\hypersetup{colorlinks=true,linkcolor=blue!50!black,urlcolor=blue!50!black}

\setlist{itemsep=1pt,topsep=3pt,parsep=0pt}
\setlength{\parindent}{0pt}
\setlength{\parskip}{3pt}
\emergencystretch=3em

% ---------- Identification de la fiche -------------------------
\newcommand{\ficheid}{}
\newcommand{\fichetitre}{}
\newcommand{\fiche}[3]{%
  \renewcommand{\ficheid}{#1}\renewcommand{\fichetitre}{#2}%
  \pagestyle{fancy}\fancyhf{}%
  \fancyhead[L]{\small\textsf{Advanced Probabilities -- Fiche #1}}%
  \fancyhead[R]{\small\textsf{#2}}%
  \fancyfoot[C]{\small\thepage}%
  \begin{center}
    {\LARGE\bfseries Fiche #1 -- #2}\par\smallskip
    {\large\itshape (#3)}\par\smallskip
    {\small Cours de V.\,M.~Lambert, M1 -- semestre 1}
  \end{center}\vspace{-2pt}\hrule\vspace{6pt}}

% ---------- Environnements numérotés (un seul compteur) --------
\theoremstyle{definition}
\newtheorem{definition}{Définition}
\renewcommand{\thedefinition}{\ficheid.\arabic{definition}}
\newtheorem{theoreme}[definition]{Théorème}
\newtheorem{proposition}[definition]{Proposition}
\newtheorem{lemme}[definition]{Lemme}
\newtheorem{corollaire}[definition]{Corollaire}
\newtheorem{exemple}[definition]{Exemple}
\theoremstyle{remark}
\newtheorem{remarque}[definition]{Remarque}

% Résultat hors cours (standard, signalé comme tel)
\theoremstyle{definition}
\newtheorem{horscours}[definition]{Hors cours}

% Idée de preuve
\newenvironment{idee}{\par\small\textit{Idée de preuve.}~}{\par}

% ---------- Encadrés -----------------------------------------
\tcbset{fichebox/.style={enhanced,breakable,boxrule=0.6pt,arc=2pt,
  left=5pt,right=5pt,top=3pt,bottom=3pt,fonttitle=\bfseries}}
\newtcolorbox{piege}[1][]{fichebox,colback=red!4,colframe=red!60!black,
  title={Pièges classiques},#1}
\newtcolorbox{liens}[1][]{fichebox,colback=blue!3,colframe=blue!50!black,
  title={Liens entre notions},#1}
\newtcolorbox{methode}[1][Méthode]{fichebox,colback=green!4,colframe=green!40!black,title={#1}}
\newtcolorbox{alerte}[1][]{fichebox,colback=orange!6,colframe=orange!70!black,
  title={Attention à l'énoncé},#1}

% Exercice de TD (fiches de rappels)
\newcounter{exo}
\newenvironment{exo}[1][]{\refstepcounter{exo}\par\medskip
  \noindent{\large\bfseries Exercice~\theexo}\ifx\relax#1\relax\else~{\large\bfseries -- #1}\fi\par\nopagebreak\smallskip}{\par}
\newcommand{\outils}{\textbf{Outils du cours.}~}
\newcommand{\meth}{\textbf{Méthode / astuce.}~}
\newcommand{\indic}{\textbf{Indication.}~}
\newcommand{\rep}{\textbf{Réponse.}~}

% ---------- Notations du cours ---------------------------------
\newcommand{\R}{\mathbb{R}}
\newcommand{\N}{\mathbb{N}}
\newcommand{\Q}{\mathbb{Q}}
\newcommand{\C}{\mathbb{C}}
\newcommand{\PP}{\mathbb{P}}
\newcommand{\QQ}{\mathbb{Q}}
\newcommand{\E}{\mathbb{E}}
\newcommand{\F}{\mathcal{F}}
\newcommand{\G}{\mathcal{G}}
\newcommand{\Hc}{\mathcal{H}}
\newcommand{\Tc}{\mathcal{T}}
\newcommand{\Bc}{\mathcal{B}}
\newcommand{\Cc}{\mathcal{C}}
\newcommand{\Ec}{\mathcal{E}}
\newcommand{\Pc}{\mathcal{P}}
\newcommand{\Nc}{\mathcal{N}}
\newcommand{\LN}{\mathcal{LN}}
\newcommand{\Expo}{\mathcal{E}}
\newcommand{\Poi}{\mathcal{P}}
\newcommand{\Ff}{\mathbb{F}}
\newcommand{\1}{\mathbf{1}}
\newcommand{\dd}{\,\mathrm{d}}
\DeclareMathOperator{\Var}{\mathbb{V}}
\DeclareMathOperator{\Cov}{Cov}
\DeclareMathOperator{\card}{card}
\DeclareMathOperator{\diag}{diag}
\newcommand{\indep}{\perp\!\!\!\perp}
\newcommand{\ps}{\ \PP\text{-p.s.}}
\newcommand{\tendloi}{\xrightarrow[n\to+\infty]{\text{loi}}}
\newcommand{\tendps}{\xrightarrow[n\to+\infty]{\text{p.s.}}}
\newcommand{\tendproba}{\xrightarrow[n\to+\infty]{\PP}}
\newcommand{\tendLp}{\xrightarrow[n\to+\infty]{L^p}}
\newcommand{\tend}{\xrightarrow[n\to+\infty]{}}

%     \fiche{C5}{Vecteurs gaussiens}{Gaussian vectors}
% ==============================================================
\usepackage[T1]{fontenc}
\usepackage[utf8]{inputenc}
\usepackage[french]{babel}
\usepackage{lmodern}
\usepackage[a4paper,margin=1.9cm,top=1.7cm,bottom=1.9cm]{geometry}
\usepackage{amsmath,amssymb,amsthm,mathtools}
\usepackage[inline]{enumitem}
\usepackage{booktabs,array,tabularx}
\usepackage{xcolor}
\usepackage[most]{tcolorbox}
\usepackage{fancyhdr}
\usepackage{hyperref}
\hypersetup{colorlinks=true,linkcolor=blue!50!black,urlcolor=blue!50!black}

\setlist{itemsep=1pt,topsep=3pt,parsep=0pt}
\setlength{\parindent}{0pt}
\setlength{\parskip}{3pt}
\emergencystretch=3em

% ---------- Identification de la fiche -------------------------
\newcommand{\ficheid}{}
\newcommand{\fichetitre}{}
\newcommand{\fiche}[3]{%
  \renewcommand{\ficheid}{#1}\renewcommand{\fichetitre}{#2}%
  \pagestyle{fancy}\fancyhf{}%
  \fancyhead[L]{\small\textsf{Advanced Probabilities -- Fiche #1}}%
  \fancyhead[R]{\small\textsf{#2}}%
  \fancyfoot[C]{\small\thepage}%
  \begin{center}
    {\LARGE\bfseries Fiche #1 -- #2}\par\smallskip
    {\large\itshape (#3)}\par\smallskip
    {\small Cours de V.\,M.~Lambert, M1 -- semestre 1}
  \end{center}\vspace{-2pt}\hrule\vspace{6pt}}

% ---------- Environnements numérotés (un seul compteur) --------
\theoremstyle{definition}
\newtheorem{definition}{Définition}
\renewcommand{\thedefinition}{\ficheid.\arabic{definition}}
\newtheorem{theoreme}[definition]{Théorème}
\newtheorem{proposition}[definition]{Proposition}
\newtheorem{lemme}[definition]{Lemme}
\newtheorem{corollaire}[definition]{Corollaire}
\newtheorem{exemple}[definition]{Exemple}
\theoremstyle{remark}
\newtheorem{remarque}[definition]{Remarque}

% Résultat hors cours (standard, signalé comme tel)
\theoremstyle{definition}
\newtheorem{horscours}[definition]{Hors cours}

% Idée de preuve
\newenvironment{idee}{\par\small\textit{Idée de preuve.}~}{\par}

% ---------- Encadrés -----------------------------------------
\tcbset{fichebox/.style={enhanced,breakable,boxrule=0.6pt,arc=2pt,
  left=5pt,right=5pt,top=3pt,bottom=3pt,fonttitle=\bfseries}}
\newtcolorbox{piege}[1][]{fichebox,colback=red!4,colframe=red!60!black,
  title={Pièges classiques},#1}
\newtcolorbox{liens}[1][]{fichebox,colback=blue!3,colframe=blue!50!black,
  title={Liens entre notions},#1}
\newtcolorbox{methode}[1][Méthode]{fichebox,colback=green!4,colframe=green!40!black,title={#1}}
\newtcolorbox{alerte}[1][]{fichebox,colback=orange!6,colframe=orange!70!black,
  title={Attention à l'énoncé},#1}

% Exercice de TD (fiches de rappels)
\newcounter{exo}
\newenvironment{exo}[1][]{\refstepcounter{exo}\par\medskip
  \noindent{\large\bfseries Exercice~\theexo}\ifx\relax#1\relax\else~{\large\bfseries -- #1}\fi\par\nopagebreak\smallskip}{\par}
\newcommand{\outils}{\textbf{Outils du cours.}~}
\newcommand{\meth}{\textbf{Méthode / astuce.}~}
\newcommand{\indic}{\textbf{Indication.}~}
\newcommand{\rep}{\textbf{Réponse.}~}

% ---------- Notations du cours ---------------------------------
\newcommand{\R}{\mathbb{R}}
\newcommand{\N}{\mathbb{N}}
\newcommand{\Q}{\mathbb{Q}}
\newcommand{\C}{\mathbb{C}}
\newcommand{\PP}{\mathbb{P}}
\newcommand{\QQ}{\mathbb{Q}}
\newcommand{\E}{\mathbb{E}}
\newcommand{\F}{\mathcal{F}}
\newcommand{\G}{\mathcal{G}}
\newcommand{\Hc}{\mathcal{H}}
\newcommand{\Tc}{\mathcal{T}}
\newcommand{\Bc}{\mathcal{B}}
\newcommand{\Cc}{\mathcal{C}}
\newcommand{\Ec}{\mathcal{E}}
\newcommand{\Pc}{\mathcal{P}}
\newcommand{\Nc}{\mathcal{N}}
\newcommand{\LN}{\mathcal{LN}}
\newcommand{\Expo}{\mathcal{E}}
\newcommand{\Poi}{\mathcal{P}}
\newcommand{\Ff}{\mathbb{F}}
\newcommand{\1}{\mathbf{1}}
\newcommand{\dd}{\,\mathrm{d}}
\DeclareMathOperator{\Var}{\mathbb{V}}
\DeclareMathOperator{\Cov}{Cov}
\DeclareMathOperator{\card}{card}
\DeclareMathOperator{\diag}{diag}
\newcommand{\indep}{\perp\!\!\!\perp}
\newcommand{\ps}{\ \PP\text{-p.s.}}
\newcommand{\tendloi}{\xrightarrow[n\to+\infty]{\text{loi}}}
\newcommand{\tendps}{\xrightarrow[n\to+\infty]{\text{p.s.}}}
\newcommand{\tendproba}{\xrightarrow[n\to+\infty]{\PP}}
\newcommand{\tendLp}{\xrightarrow[n\to+\infty]{L^p}}
\newcommand{\tend}{\xrightarrow[n\to+\infty]{}}

\begin{document}
\fiche{C9}{Processus de Poisson}{Poisson process}

\textit{Source : slides, chapitre 10 (Thm~21) ; poly « Gaussian vector » (Déf.~33, Thm~34--39), plus détaillé.
Modèle d'arrivées aléatoires : clients d'une file, appels, sinistres (risque), blocs d'une blockchain.}

\section*{1. Processus ponctuel et fonction de comptage}

\begin{definition}[processus ponctuel, \emph{point process}]\label{C9:ponctuel}
Sur $(\Omega,\F,\PP)$, une suite $(T_n)_{n\ge1}$ de v.a.\ à valeurs dans $\R_+$, strictement croissante, avec $T_n\to+\infty$ p.s. ;
$T_0=0$ par convention : $0<T_1<T_2<\dots<T_n<\dots\to\infty$. $T_n$ est l'instant de la $n$-ième arrivée.
Inter-arrivées : $\tau_n=T_n-T_{n-1}>0$, et $T_n=\sum_{i=1}^n\tau_i$ (connaître $(T_n)$ ou $(\tau_n)$, c'est pareil).
\end{definition}

\begin{definition}[fonction de comptage, \emph{counting process}]\label{C9:comptage}
$N_t=\sup\{j\ge1,\ T_j\le t\}=\sum_{j\ge1}\1_{\{T_j\le t\}}$ ($=0$ si $T_1>t$) : nombre d'arrivées dans $[0,t]$.
$N_0=0$ p.s.\ et $N_t<\infty$ p.s.\ (car $T_n\to\infty$).
\end{definition}

\begin{proposition}[dictionnaire $N\leftrightarrow T$]\label{C9:dico}
Pour $n\ge1$ et $0\le s<t$ :
\begin{gather*}\{N_t\ge n\}=\{T_n\le t\},\qquad \{N_t=n\}=\{T_n\le t<T_{n+1}\},\\
\{N_t\le n\}=\{T_{n+1}>t\},\qquad \{N_s<n\le N_t\}=\{s<T_n\le t\}.\end{gather*}
\end{proposition}
\begin{idee}
Si $T_n(\omega)\le t$, comme $T_1\le\dots\le T_n$, on a $T_j(\omega)\le t$ pour $j\le n$, donc $N_t(\omega)\ge n$ ; réciproquement,
$N_t(\omega)\ge n$ force $T_j(\omega)\le t$ pour $j\le n$. Le reste par complémentaire :
$\{N_t\le n\}=\{N_t\ge n+1\}^c=\{T_{n+1}\le t\}^c$, et $\{N_t=n\}=\{N_t\ge n\}\cap\{N_t\le n\}$.
\end{idee}

\section*{2. Processus de Poisson}

\begin{definition}[processus de Poisson d'intensité $\lambda>0$, slides]\label{C9:def}
Un processus de comptage $N=(N_t)_{t\ge0}$ tel que
\begin{enumerate}
  \item $N_0=0$ ;
  \item \textbf{accroissements indépendants} : pour $0\le t_1<\dots<t_n$, $N_{t_1},N_{t_2}-N_{t_1},\dots,N_{t_n}-N_{t_{n-1}}$ sont indépendantes ;
  \item \textbf{accroissements stationnaires} : pour $s,t\ge0$, $N_{t+s}-N_s$ a même loi que $N_t$ ;
  \item pour tout $t\ge0$, $N_t\sim\Poi(\lambda t)$.
\end{enumerate}
C'est un processus de Lévy (accroissements indépendants et stationnaires).
\end{definition}

\begin{theoreme}[poly : les accroissements suffisent]\label{C9:thm34}
Si $N$ est un processus de comptage à accroissements indépendants et stationnaires (Déf.~33 du poly), il existe $\lambda>0$ tel que
pour $0\le s<t$ : $N_t-N_s\sim\Poi(\lambda(t-s))$, i.e.\ $\PP(N_t-N_s=k)=e^{-\lambda(t-s)}\frac{(\lambda(t-s))^k}{k!}$.
En particulier $N_t\sim\Poi(\lambda t)$, et $\lambda=\E(N_{t+1}-N_t)$ est le nombre moyen d'arrivées par unité de temps.
\end{theoreme}
\begin{idee}
$G_t(u)=\E(u^{N_t})$ vérifie, par indépendance puis stationnarité, $G_t(u)=\E(u^{N_t-N_s})\E(u^{N_s})=G_{t-s}(u)G_s(u)$,
d'où $G_t(u)=(G_1(u))^t=e^{-t\lambda(u)}$ ; il reste à montrer $\lambda(u)=\lambda(0)(1-u)$, fonction génératrice de Poisson.
\end{idee}

\begin{theoreme}[première arrivée et attente résiduelle]\label{C9:T1}
$T_1\sim\Expo(\lambda)$. Pour tout $s\ge0$, l'attente après $s$ de la première arrivée, $T_{N_s+1}-s$, suit aussi $\Expo(\lambda)$.
\end{theoreme}
\begin{idee}
$\{T_1>t\}=\{N_t\ge1\}^c=\{N_t=0\}$ donc $\PP(T_1>t)=e^{-\lambda t}$, la queue de $\Expo(\lambda)$.
$\{T_{N_s+1}-s>t\}=\{N_{s+t}-N_s=0\}$ (aucune arrivée dans $]s,s+t]$), de probabilité $\PP(N_t=0)=e^{-\lambda t}$ par stationnarité.
\end{idee}

\begin{theoreme}[propriété de Markov forte, poly Thm~36]\label{C9:markov}
Soit $S$ un $\Ff^N$-temps d'arrêt avec $S<\infty$ p.s., et $N^S_t=N_{S+t}-N_S$. Alors $(N^S_t)_{t\ge0}$ est un processus de Poisson
d'intensité $\lambda$, indépendant de $\F^N_S$. ($N^S$ compte les arrivées $T^S_n=T_{N_S+n}-S$ après $S$.)
\end{theoreme}
% À VÉRIFIER : le poly écrit « Assume that {S < infinity} » (lu : S < infini p.s.) et « Then (N_t) is a Poisson process » (lu : N^S).
\begin{remarque}
\emph{Hors cours :} $\F^N_S=\{A\in\F,\ A\cap\{S\le t\}\in\F^N_t\ \forall t\}$ est l'information disponible à l'instant aléatoire $S$.
\end{remarque}

\begin{theoreme}[inter-arrivées, slides Thm~21 / poly Thm~37]\label{C9:interarrivees}
Si $N$ est un processus de Poisson d'intensité $\lambda$, les $(\tau_n)_{n\ge1}$ sont \textbf{indépendantes et identiquement distribuées}
de loi $\Expo(\lambda)$.
\end{theoreme}
\begin{idee}
$\tau_1=T_1\sim\Expo(\lambda)$ (Thm~\ref{C9:T1}). Markov forte avec $S=T_1$ : $N^{T_1}$ est un Poisson($\lambda$) indépendant de
$\F^N_{T_1}\ni\tau_1$, et sa première arrivée est $T_2-T_1=\tau_2$. On itère.
\end{idee}

\begin{theoreme}[réciproque, poly Thm~38]\label{C9:construction}
Si $(\tau_n)_{n\ge1}$ sont indépendantes de loi $\Expo(\lambda)$, $T_0=0$, $T_n=\sum_{i=1}^n\tau_i$ et $N_t=\sup\{j\ge1,\ T_j\le t\}$, alors
$(N_t)_{t\ge0}$ est un processus de Poisson d'intensité $\lambda$.
\end{theoreme}

\begin{theoreme}[comportement en temps long, poly Thm~39]\label{C9:asympt}
$\dfrac{N_t}{t}\xrightarrow[t\to+\infty]{\text{p.s.}}\lambda$ \quad et \quad
$\dfrac{N_t-\lambda t}{\sqrt{\lambda t}}\xrightarrow[t\to+\infty]{\text{loi}}\Nc(0,1)$.
\end{theoreme}
\begin{idee}
$N_n=\sum_{i=1}^n(N_i-N_{i-1})$, somme de v.a.\ i.i.d.\ $\Poi(\lambda)$ : LGN puis TCL aux temps entiers (C4) ; on contrôle ensuite
$\frac{N_t}{t}=\frac{N_{[t]}}{[t]}\cdot\frac{[t]}{t}+\frac{N_t-N_{[t]}}{t}$ avec $\sup_{n<t<n+1}\frac{N_t-N_n}{n}\to0$ (« encore un peu de travail »).
\end{idee}

\begin{horscours}[à savoir, standard]\label{C9:hc}
$\E(N_t)=\Var(N_t)=\lambda t$. $T_n$ (somme de $n$ exponentielles indépendantes) suit la loi d'Erlang / Gamma$(n,\lambda)$ (densité rappelée dans les TD4 et TD5) :
$f_{T_n}(t)=\lambda^n\frac{t^{n-1}}{(n-1)!}e^{-\lambda t}\1_{t>0}$. Si $X\sim\Poi(\lambda)\indep Y\sim\Poi(\mu)$, $X+Y\sim\Poi(\lambda+\mu)$.
Absence de mémoire de l'exponentielle : $\PP(\tau>s+t\mid\tau>s)=\PP(\tau>t)$.
\end{horscours}

\begin{liens}
\[\text{Déf.~\ref{C9:def} (accroissements + loi de Poisson)}\iff\text{inter-arrivées i.i.d.\ }\Expo(\lambda)\quad(\text{Thm~\ref{C9:interarrivees} et \ref{C9:construction}})\]
\begin{itemize}
  \item Accroissements indépendants et accroissements stationnaires sont deux hypothèses distinctes : aucune n'entraîne l'autre.
  Par exemple, un processus de Poisson d'intensité variable $\lambda(t)$ (hors cours) garde des accroissements indépendants mais
  non stationnaires. Réciproquement, \emph{stationnaires $\not\Rightarrow$ indépendants} (hors cours) : processus de Poisson \emph{mélangé}, où $\Lambda$ vaut 1 ou 2 avec probabilité $\frac12$ et, sachant $\Lambda$, $N$ est un Poisson d'intensité $\Lambda$. Ses accroissements sont stationnaires, mais $\Cov(N_1,N_2-N_1)=\Var(\Lambda)=\frac14>0$.
  \item Poisson $\Rightarrow$ $N_t\sim\Poi(\lambda t)$ pour tout $t$ ; mais les lois marginales $(N_t)$ seules ne déterminent pas la loi du processus : il faut les lois fini-dimensionnelles (C8) ou les hypothèses sur les accroissements.
\end{itemize}
\end{liens}

\begin{piege}
\begin{itemize}
  \item $\Expo(\lambda)$ a pour moyenne $1/\lambda$ ; $\Poi(\lambda t)$ a pour moyenne $\lambda t$ ; ne pas confondre intensité et durée moyenne.
  \item $\{N_t\ge n\}=\{T_n\le t\}$ : le sens des inégalités s'inverse entre $N$ et $T$.
  \item L'intervalle $T_{N_t+1}-T_{N_t}$ \emph{qui contient} $t$ n'est pas de loi $\Expo(\lambda)$ : il est en moyenne plus long
  (paradoxe de l'autobus, TD5, ex.~11). En revanche l'attente résiduelle $T_{N_t+1}-t$ est bien $\Expo(\lambda)$.
  \item Pour calculer une probabilité sur $T_n$, penser à la traduire en événement sur $N_t$ (et réciproquement).
  \item Les accroissements doivent porter sur des intervalles \emph{disjoints} pour être indépendants.
\end{itemize}
\end{piege}

\section*{Exemples types}

\begin{exemple}[calcul via le dictionnaire]\label{C9:ex1}
$\lambda=1{,}5$ par heure, $t=2$ h, donc $\lambda t=3$. $\PP(N_2=3)=e^{-3}\frac{3^3}{3!}=4{,}5\,e^{-3}\approx0{,}224$.
$\PP(T_3\le2)=\PP(N_2\ge3)=1-e^{-3}\bigl(1+3+\frac92\bigr)=1-8{,}5\,e^{-3}\approx0{,}577$.
\end{exemple}

\begin{exemple}[attente résiduelle]\label{C9:ex2}
Clients à $\lambda=2$ par minute. Quel que soit le passé (indépendance vis-à-vis de $\F^N_s$, Thm~\ref{C9:markov} avec $S=s$), la probabilité qu'aucun client n'arrive dans les 30 prochaines secondes est
$\PP(T_{N_s+1}-s>\frac12)=e^{-2\cdot\frac12}=e^{-1}\approx0{,}368$ (Thm~\ref{C9:T1}), et l'attente moyenne est $\frac1\lambda=30$ s.
\end{exemple}

\end{document}
